\section{Derivación del Attention Score}

\subsection{Expansión del Pre-Softmax Score}

Bajo RoPE, el attention score (pre-softmax) entre la Query de la posición $m$ y la Key de la posición $n$ es:

$$
s(m,n) = (R_m q)^T (R_n k) = q^T R_m^T R_n k = q^T R_{n-m} k
$$

Al expandir $R_{n-m}$ en cada subespacio 2D, la contribución de cada subespacio $j$ es proporcional a:

$$
\cos\big((m-n) \cdot \theta_j\big)
$$

\begin{importantbox}{Fórmula central}
El attention score es proporcional a la suma de las funciones coseno de todos los subespacios:
$$
s(m,n) \propto \sum_{j=0}^{d/2-1} \cos\big((m-n) \cdot \theta_j\big)
$$
Cada subespacio aporta una onda coseno de una frecuencia distinta, determinada por $\theta_j$.
\end{importantbox}

\subsection{El efecto de la distancia $\tau = m - n$}

Sea $\tau = m - n$ la distancia entre tokens; analicemos el comportamiento de $s(\tau)$:

\textbf{Distancia corta ($\tau$ pequeña)}:
\begin{itemize}
  \item En todos los subespacios, $\cos(\tau \cdot \theta_j) \approx 1$
  \item Los términos se suman con signo positivo y el score es muy alto
  \item El modelo atiende fuertemente a los tokens cercanos
\end{itemize}

\textbf{Distancia larga ($\tau$ grande)}:
\begin{itemize}
  \item Subespacios de alta frecuencia ($\theta_j$ grande, $j$ pequeña): $\tau \cdot \theta_j$ cambia bruscamente y el valor de $\cos$ oscila con rapidez
  \item Subespacios de baja frecuencia ($\theta_j$ pequeña, $j$ grande): $\tau \cdot \theta_j$ cambia lentamente y el valor de $\cos$ decrece con lentitud
  \item Las oscilaciones bruscas de la parte de alta frecuencia se \textbf{cancelan entre positivos y negativos} al sumarse
\end{itemize}

\subsection{La imagen física de la atenuación remota}

\begin{importantbox}{El mecanismo de atenuación remota de RoPE}
La atenuación a larga distancia de RoPE es el resultado conjunto de dos efectos:
\begin{enumerate}
  \item \textbf{Cancelación total de las dimensiones de alta frecuencia}: los subespacios del lado izquierdo rotan con rapidez y, a distancias grandes, se cancelan entre positivos y negativos (de forma similar al efecto de interferencia de ondas)
  \item \textbf{Descenso lento de las dimensiones de baja frecuencia}: los subespacios del lado derecho rotan con lentitud y conservan una asociación débil a distancia
\end{enumerate}
En conjunto: matemáticamente, RoPE \textbf{atiende principalmente a lo local}, pero \textbf{conserva la capacidad remota}.
\end{importantbox}

Esto se asemeja al \textbf{fenómeno de interferencia de ondas} de la física: ondas de distintas frecuencias se cancelan entre sí a distancia, pero ciertas frecuencias específicas (las bajas) siguen propagándose hasta lejos.

\subsection{Análisis de los casos extremos}

Para el subespacio más a la derecha ($j = d/2 - 1$):
$$
\theta_{d/2-1} = \text{base}^{-1} \approx \frac{1}{\text{base}}
$$

Si base = 10000, entonces $\theta \approx 10^{-4}$, una longitud de onda extremadamente larga. Aunque $\tau$ sea muy grande, el crecimiento de $\tau \cdot \theta$ sigue siendo muy lento y $\cos(\tau \cdot \theta) \approx 1$. Esta es la razón por la que se conserva la información remota.

\subsection{Resumen de la sección}

El attention score es la superposición de las ondas coseno de todos los subespacios. A larga distancia, los términos de alta frecuencia se cancelan y los de baja frecuencia se atenúan lentamente; juntos producen la propiedad de atenuación remota de RoPE. RoPE implementa de manera natural un patrón de atención de «atender a lo local + conservar lo remoto».

